This work presents the development of a web-based electronic health record system for the Dental Rehabilitation and TMD Academic League of the State University of Ponta Grossa (LORD-UEPG). The motivation lies in the limitations of paper records in the follow-up of long-term treatments, such as those for temporomandibular disorders (TMD) and prosthetic rehabilitation, in which the information recorded at each visit must be compared over time. The system digitizes the League's clinical form from a versioned data schema, stored in a hybrid relational and document model (PostgreSQL with JSONB), with role-based access control, faculty supervision workflow, audit trail and pseudonymization of research data, in compliance with the Brazilian General Data Protection Law. It includes a locally executed artificial intelligence module that transcribes the student's speech and organizes the anamnesis report into the form fields without issuing any diagnosis: each suggestion is reviewed individually by a human. An analytical module in Python computes clinical progress indicators, such as pain scales, the RDC/TMD Graded Chronic Pain Scale and mouth opening. On 30 synthetic clinical cases, the local 3-billion-parameter language model reached an F1 score of 0.59, and deterministic safeguards reduced from 32\% to 25\% the proportion of empty fields wrongly filled, which confirms the need for human review. Speech transcription, using a synthetic voice, missed about 4\% of the words in silence and 6\% with simulated noise, with errors concentrated in clinical terms. With 10,000 patients, the hybrid model read a patient record three to four times faster than a normalized model and used almost half the storage. The usability evaluation, carried out without participants, found no accessibility violations in the main screens and estimated about 9.7 minutes to fill in the complete form. On a synthetic cohort with known effects, the analytical module recovered the pain reductions programmed for the three treatments. The results indicate that it is feasible to keep the data within the institution using free tools and to use AI only to support clinical documentation.

\vspace{\onelineskip}

\noindent\textbf{Keywords}: Electronic health records. Dentistry. Temporomandibular disorders. Artificial intelligence. Data protection.
